%!TEX program = lualatex
\documentclass[ngerman,headheight=70pt]{scrartcl}
\usepackage[ngerman]{babel}
\usepackage{blindtext}
\usepackage[utf8]{luainputenc}
\usepackage{scrlayer-scrpage}
\usepackage[T1]{fontenc}
\usepackage{fontspec}
\usepackage{multicol}
\usepackage{enumitem}
\usepackage{epigraph}
\usepackage[german=quotes]{csquotes}
\usepackage[bottom=2cm,footskip=8mm]{geometry}
\setmainfont{Source Sans Pro}
\MakeOuterQuote{"}

\setlength{\parindent}{0pt}
\setlength{\parskip}{6pt}

\pagestyle{plain.scrheadings}
\thispagestyle{scrheadings}
\lohead{\vspace*{-1.5cm}\includegraphics[scale=0.4]{../images/stupa_logo_rot.pdf}}
\rohead{\\Präsidium des\\Studierendenparlaments}

\lofoot[Universität Hamburg · Präsidium des Studierendenparlaments\\\hrule Von-Melle-Park 5 · D-20146 Hamburg]{Universität Hamburg · Präsidium des Studierendenparlaments\\\hrule Von-Melle-Park 5 · D-20146 Hamburg}
\cofoot[]{}
\cohead[\thepage]{}

\renewcommand{\thesubsection}{TOP \arabic{subsection}}
\renewcommand{\thesubsubsection}{\arabic{subsubsection}.)}
\setcounter{subsection}{1}

\begin{document}
    UHH · StuPa-Präsidium · Von-Melle-Park 5 · D-20146 Hamburg

    \section*{Beschluss des Studierendenparlamentes vom 28. April 2016}
    \subsection*{Hochschul- und Wissenschaftspolitik als Gesellschaftspolitik}

	\subsubsection{Alles drängt nach einer Abkehr vom Neoliberalismus.}

    Die prinzipielle Krisenhaftigkeit und darin die aktuelle, tiefe Krise des
    Kapitalismus offenbart mehr und mehr dessen Unfähigkeit, selbst die
    grundlegenden Bedürfnisse und Ansprüche auf ein "gutes Leben für Alle" zu
    erfüllen. Der Status quo -- meint die Politik aus wettbewerblichem Umbau des
    Sozialstaats, Privatisierung der öffentlichen Infrastruktur und
    Standortkonkurrenz zur Sicherung der Profite Weniger -- wird zunehmend in
    Frage gestellt. Die neoliberale Erzählung, dass der Markt es schon regelt,
    überzeugt vor dem Hintergrund der sich rasant zuspitzenden sozialen
    Ungleichheit, der zunehmenden Ökonomisierung aller Lebensbereiche und den
    kriegerischen Konflikten in der Welt immer weniger Menschen.

    Das politische, neoliberale Establishment verliert massiv den Boden unter
    den Füßen: während die soziale Alternative zum krisenhaften Status quo durch
    linke und solidarische Bewegungen zunehmend mutig und massenhaft vertreten
    wird, zeigt die immer brutaler auftretende gesellschaftliche Rechte in
    Richtung Barbarei. Es drängt sich jedem auf, dass eine Entscheidung, welche
    Tendenz die gesellschaftliche Entwicklung nehmen wird, ansteht.

    Diese Polarisierung zeigt sich auf der einen Seite beispielsweise in der
    linken, solidarischen Krisenbeantwortung, die in Südeuropa dominiert oder an
    der sogar in den imperialistischen Staaten USA und UK zunehmenden Popularität
    der beiden linken Massenbewegungen um Bernie Sanders und Jeremy Corbyn. Die
    vielerorts weit entwickelten Ansprüche an eine offene, soziale, demokratische
    und ökologische Gesellschaft verschaffen sich zunehmend Ausdruck u.a. mit dem
    klaren NEIN zum Kommerzwahnsinn Olympia und dem Protest gegen das
    Freihandelsabkommen TTIP, der im vergangenen Oktober mit 250.000 Menschen auf
    Berlins Straßen kulminierte. Die zivilgesellschaftliche Solidarität zur
    umfassenden gesellschaftlichen Teilhabe geflüchteter Menschen zeugt von der
    verbreiteten Haltung, dass die einzige humane Perspektive die gemeinsame der
    Unterdrückten ist.

    Gegen das alltägliche Engagement für eine solidarische, die sozialen
    Grundlagen kritisch reflektierende und verändernde Beantwortung der Krise
    steht die zunehmend skrupelloser auftretende gesellschaftliche Rechte. Diese
    offenbart ihren menschenfeindlichen Kern in massiv gesteigerter
    rassistischer Hetze und Angriffen gegen Geflüchtete. Gerade in Ländern, wo
    die Linke defensiv oder durch Opportunismus besonders unglaubwürdig ist,
    sind offen sozialdarwinistische Einstellungen zum Teil hegemonial -- der
    Blick nach Ungarn, Polen und Frankreich zeigt (graduell verschieden) wie
    ernst die Gefahr auf der rechten Seite der politischen Auseinandersetzung
    ist. Indem die Rechte Konkurrenz und soziale Ungleichheit forciert und sich
    dabei als Alternative gibt, offenbart sie zunehmend auch ihren systemwahrenden
    Gehalt, insbesondere indem nationalistische Ideologie die sozialen
    Antagonismen zwischen Arbeit und Kapital verschleiern sollen.

    Es gilt, dass wir uns in einer zugespitzten Entscheidungssituation über die
    Tendenz gesellschaftlicher Entwicklung befinden: es kommt darauf an, die
    progressive Dynamik zu echten Durchbrüchen zu bringen! Damit der Bruch mit
    dem Neoliberalismus positiv gelingen kann, kommt es zentral darauf an, Alle
    zu ermuntern, ihre sozial fortschrittlichen Ambitionen zur Geltung zu bringen
    und philosophisch-geschichtsbewusst zu fundieren. Das ist auch die
    entscheidende Gegnerschaft gegen Rechts.

    \subsubsection{Die neoliberale Rede, der Mensch sei eine Ich-AG, war von
    Anfang an eine Lüge, denn der Mensch ist ein soziales Wesen. Das zeigt sich
    täglich im Engagement von vielen:}

    Ökonomisierung für mehr Personal streikenden GesundheitspflegerInnen der
    Berliner Charité, ob die Drittmittel und LOMI ablehnenden ForscherInnen am
    Soziologie-Institut der Uni Bielefeld, in der auf "Bildung mündiger
    Persönlichkeiten" zielenden Studienreform der UHH oder die millionenfach in
    der Flüchtlingssolidarität Aktiven, es wird bereits sehr weitgehend eine
    andere Alltagskultur praktiziert, als es die neoliberalen FürsprecherInnen
    und Strukturen gerne hätten.

    Die organisiert progressiven Kräfte haben die Verantwortung, diese --
    teilweise noch unbewusste -- kulturelle Alternative zu politisieren und Alle
    zu ermutigen, hin zu einer verändernden Praxis.  Dabei gilt es vor allem,
    sich aus der neoliberalen Ideologie der Unbedeutendheit und
    Ungesellschaftlichkeit des Subjekts ("there is no such thing as society.
    There are individual men and women, and there are families"; Margaret
    Thatcher) herauszuarbeiten. Kollektive -- und darin subjektive --
    Handlungsfähigkeit erlangen wir nur in bewusster Gegnerschaft zu dieser
    Einschränkung und der auf Einschüchterung progressiver Ansprüche zielender
    Hetze der Rechten. Wir müssen jeden Tag mit jeder Facette unseres Wirkens
    eine Kultur schaffen, die den radikalen gesellschaftlichen Veränderungswillen
    als einzig realistische Position bestärkt und uns so gegen alle Kräfte der
    Hoffnungslosigkeit, des Zynismus und des Expertentums wenden.

    \blockquote["Gefängnishefte", Heft 10, Teil II, § (54) -- Antonio Gramsci, 1932]{
    Denn der einzelne kann sich mit all denen zusammenschließen, die dieselbe
    Veränderung wollen, und wenn diese Veränderung vernünftig ist, kann der
    Einzelne sich in einem imponierenden Ausmaß vervielfachen und eine
    Veränderung erzielen, die viel radikaler ist, als es auf den ersten Blick
    möglich erscheint.}

    Deutlich wurde dies auch an der -- von der Studierendenschaft maßgeblich
    wissenschaftlich und auf die positive Alternative zielend gestalteten --
    NOlympia-Kampagne in Hamburg, die der Skepsis gegenüber einer
    Stadtentwicklung nach dem Prinzip HafenCity und dem Engagement für eine
    soziale, demokratische und ökologische Stadt in einem kraftvollen NEIN beim
    Referendum Ausdruck verlieh.

    \subsubsection{Dagegen ist ideologisch und praktisch die zentrale Waffe der
    fortgeführten Kommerzialisierung aller Lebensbereiche nach wie vor die
    "Schuldenbremse".}

    Hauptfunktion ist Einschüchterung und Unterdrückung progressiver
    gesellschaftlicher Entwicklungsansprüche. Scharfe Offensivität gegen
    Barbarei, Perspektivlosigkeit und Irrealität der Austeritätspolitik sowie
    Hoffnung und Engagement für eine menschenwürdige Entwicklungsperspektive
    bilden eine Einheit.

    \subsubsection{In der Bundesrepublik ist die Abkehr von der Unterwerfung
    unter das Prinzip der Gewinnmaximierung in der Hochschul- und
    Wissenschaftspolitik sehr weit verallgemeinerungswürdig ausgeprägt.}

    Dies ist ganz wesentlich dem Umstand geschuldet, dass bereits die europaweite
    studentische Protestbewegung 2009, die ihren Ausgang in der Gegnerschaft zum
    "Humankapital"-Prinzip der Bologna-Reform nahm, über einen Abwehrkampf
    hinaus ging und für die positive Perspektive der "Bildung mündiger Menschen"
    kämpfte. Die Bedingungen für die Entwicklung gesellschaftlicher Opposition
    sind in Hochschulen erkämpfter Weise verhältnismäßig gut, da Wissenschaft dem
    Kern nach argumentativ-demokratisch und wahrheitssuchend ist. Ihre Mitglieder
    haben die gesellschaftliche Verantwortung und die Möglichkeit, als
    wissenschaftliche Institution solidarische Lösungsansätze für
    gesellschaftliche Herausforderungen zu erarbeiten, dafür hegemoniebildend
    einzugreifen und so zur humanen Lösung der Krise beizutragen. Deswegen ist
    die Hochschule ein umkämpftes Feld zwischen von den Herrschenden forcierter
    Reproduktion des Bestehenden und emanzipatorischer Kritik sowie Veränderung
    gesellschaftlicher Verhältnisse. Der Sinn der progressiven Studienreform
    besteht darin, die wissenschaftliche Aneignung der Welt zu ihrer kollektiven
    menschenwürdigen Gestaltung zu verallgemeinern.

    \subsubsection{Mit der Exzellenzinitiative soll diese Tendenz abgelenkt und
    eingeschränkt werden.}

    Deshalb ist für AStA und VS die Ermutigung, dazu Nein zu sagen, eine
    wesentliche Kernaufgabe. Sie besteht auch darin, die so gefassten Ansprüche
    voll zur Geltung zu bringen. Auf der Grundlage jahrzehntelanger bewusster
    Unterfinanzierung sollen die Hochschulen mit der Exzellenzinitiative weiter
    gegeneinander und in sich in Konkurrenz gesetzt werden. So soll zur
    Konformität gezwungen werden. Nur noch Wenige halten die Exzellenzinitiative
    für wissenschaftlich sinnvoll. Positive Durchbrüche in der Hochschul- und
    Wissenschaftspolitik sind gesellschaftlich beispielgebend für die Auflösung
    der vermeintlichen Alternativlosigkeit.

    \subsubsection{Für die positive Wendung der Krise tragen die VS und der AStA
    eine entscheidende Verantwortung. Wir haben Initiativ-, Aufklärungs- und
    Perspektiv-Funktion.}

    Es geht um echte soziale und kulturelle Verbesserungen. Der brutalste
    Ausdruck der Überkommenheit des Kapitalismus ist der Krieg. Frieden ist
    dagegen die klarste Alternative.

    \subsubsection{Der AStA und die VS fokussieren daher ihre politische Strategie
    (weiter) darauf,}

    \begin{itemize}
        \item die Schuldenbremse zu knacken zugunsten einer auf eine
              nachfrageorientierte Investitionspolitik und Umverteilung "von oben
              nach unten" gerichteten Reforminitiative als befreiende Entwicklung
              in Sozialem, Bildung, Kultur und Gesundheit (dafür ist
              bündnispolitisch vor allem an "Etwas Besseres als Olympia"
              anzuknüpfen),
        \item für die Beendigung von Kriegen weltweit zu kämpfen zugunsten
              einer zivilen internationalen Entwicklung (Zivilklausel in Hochschule
              und Hafen),
        \item dem Leitbild der "unternehmerischen Hochschule" praktisch den
              Garaus zu machen zugunsten gesellschaftlich verantwortungsvoller
              Wissenschaften und der Bildung kritischer Subjekte (Exzellenz-Ini
              stoppen, Rahmenprüfungsordnung als Emanzipationsstruktur
              ausgestalten),
        \item die Studienreform unter dem Leitbild der Verwirklichung von Frieden
              und Menschenrechten weiter auf gesellschaftlich progressives
              Eingreifen hin voranzutreiben,
        \item die soziale Lage der Studierenden zum selbstermächtigenden
              Ausgangspunkt eines gesamtgesellschaftlichen Kampfes für eine
              emanzipatorische, unentgeltliche, barrierefreie soziale
              Infrastruktur zu machen (u.a. Rückeroberung des Studierendenwerks
              als demokratische, soziale Institution der kollektiven
              Daseinsgestaltung studentischen Lebens, Kampagne für einen
              entgeltfreien ÖPNV, Zugang zu gutem Wohnraum für Alle und
              Ausweitung von "Open Access"),
        \item anknüpfend an die reichhaltige antifaschistische Tradition von
              Arbeiterbewegung, Weißer Rose, AutorInnen verbrannter Bücher und
              "68er" und in harter Gegnerschaft zur sich radikalisierenden
              Rechten weiter die "Vernichtung des Nazismus mit seinen Wurzeln"
              und den "Aufbau einer neuen Welt des Friedens und der Freiheit"
              (Schwur von Buchenwald) zur Leitlinie all unseres Wirkens zu machen
              (tätiges Erinnern an deutschen Faschismus und Widerstand, Aufklärung
              gegen alle Ungleichheitsideologien (u.a. Sexismus,
              Altersdiskriminierung und Antikommunismus), solidarische Ermunterung
              zur Verwirklichung sozial progressiver Ambitionen entgegen
              neoliberal-kapitalistischer Passivierung). Wichtig ist dabei auch,
              über die systemerhaltende Funktion rechter Ideologie und Praxis
              aufzuklären.
        \item die ökologische Frage des nachhaltigen Austauschverhältnisses
              zwischen Mensch und Natur zum Gegenstand von Aufklärungsarbeit und
              politischer Intervention zu machen. Insbesondere die dringend
              notwendige Energiewende hin zu dezentraler, vergesellschafteter
              erneuerbarer Energiegewinnung muss wissenschaftlich weiter
              forciert werden.
        \item AStA-Kulturveranstaltungen gegen universitäre und gesellschaftliche
              Missstände zu organisieren, die zur Diskussion anregen und über die
              über Ursachen aufklärt und somit eine Alternative zu den
              kommerziellen auf reine Ablenkung gerichteten Spaßveranstaltungen
              geboten werden. Kultur ist als gestalterischer Ausdruck der geistigen,
              kritischen und praktischen Auseinandersetzung und Umformung dessen,
              was einen täglich umgibt, zu verstehen und als solche Aufklärung.
              Insbesondere über die künstlerische, filmische und literarische
              Aneignung widerständiger Geschichte, hat die VS die Aufgabe zu
              einem geschichtsbewussten Sein zu ermuntern.
        \item die soziale Öffnung der Hochschule voll zu realisieren und
              insbesondere die gesellschaftliche Teilhabe von Flüchtlingen über
              Bildung durch Wissenschaft auch sozial zu unterlegen.
    \end{itemize}

    \textbf{Die Verantwortung besteht darin die alltäglich praktizierte
    Solidarität auszubauen, hin zu einer befreienden und verändernden Kultur
    und Praxis.}

    Die Verfasste Studierendenschaft (FSRe, StuPa, AStA) muss der solidarische
    Zusammenhang im Gegensatz zu konsumistischem Turbo-Studium,
    (Eltern-)Erwartungsdruck, Burn-Out, individualisierter Prekarität und
    Drittmittel-Wissenschaft sein. So ermutigen wir aus der alltäglichen
    Konkurrenzhetze aus- und für Verbesserungen einzutreten.
\end{document}
